\subsection{Relación entre las vistas de arquitectura}
\label{sec:iso42010}

RT-02.03 exige cinco vistas de la arquitectura: lógica, procesos, despliegue, datos y seguridad. El Anexo 4-S define sus interesados, preocupaciones, modelos y reglas de correspondencia, siguiendo la organización de descripciones de ISO/IEC/IEEE 42010:2022 (ISO, 2022a). La integración atraviesa esas vistas; no se presenta como una sexta vista exigida por la Base.

Operaciones necesita comprobar que un pedido puede avanzar sin perder su lote ni saltarse una autorización; Calidad necesita seguir un lote hasta su receptor; Seguridad necesita distinguir quién puede liberar una carga; y TI necesita aislar una falla y recuperar el estado. Las vistas responden a esas preocupaciones con los mismos identificadores M1--M12, INT-01--15 y los componentes transversales del Anexo 4-N. Así se puede pasar de una secuencia a su dueño de datos, contrato y control sin cambiar de vocabulario.

Las correspondencias se comprueban en ambos sentidos:

\begin{itemize}
  \item todo evento tiene productor y consumidores;
  \item toda escritura tiene dueño;
  \item toda función desconectada tiene persistencia y autorización locales;
  \item y todo componente desplegado debe realizar una responsabilidad inventariada. Una diferencia se registra con requisito afectado y ADR responsable. La cobertura lógica del inventario no se confunde con el cotejo de emplazamientos y centros de datos.
\end{itemize}

La vista lógica desarrolla aquí las responsabilidades y sus relaciones. El despliegue y los centros de datos se acreditan en 4.2--4.3.


